\documentclass[11pt]{article}
\usepackage[T1]{fontenc}
\usepackage{lmodern}
\usepackage[margin=1in]{geometry}
\usepackage{booktabs}
\usepackage{listings}
\usepackage{amsmath}
\usepackage{float}
\usepackage{microtype}
\usepackage{xcolor}
\usepackage{setspace}
\usepackage[font=small,labelfont=bf,skip=8pt]{caption}
\usepackage{hyperref}
\setlength{\emergencystretch}{3em}
% Allow DOI/GitHub URLs to break cleanly in narrow columns.
\makeatletter
\g@addto@macro{\UrlBreaks}{\do\/\do-\do.\do\_\do:\do0\do1\do2\do3\do4\do5\do6\do7\do8\do9}
\makeatother

\hypersetup{
  colorlinks=true,
  linkcolor=black,
  citecolor=black,
  urlcolor=blue!60!black,
  pdftitle={Trail: A Deterministic-First Log Template Miner with a Per-Decision Audit Trail},
  pdfauthor={Eduard Beloiu},
  pdfsubject={Technical report on auditable log template mining},
  pdfkeywords={log parsing, log template mining, reproducibility, audit trail}
}
\lstset{
  basicstyle=\ttfamily\footnotesize,
  breaklines=true,
  breakatwhitespace=true,
  columns=fullflexible,
  keepspaces=true,
  showstringspaces=false,
  frame=single,
  aboveskip=1.4ex,
  belowskip=1.4ex,
  xleftmargin=0.4em,
  framexleftmargin=0.4em
}
\newcommand{\wild}{\texttt{<*>}}
\newcommand{\doiref}[1]{\href{https://doi.org/#1}{\texttt{#1}}}

% Slightly more air between paragraphs, floats, and section heads.
\setstretch{1.045}
\setlength{\parskip}{0.35ex plus 0.15ex}
\setlength{\floatsep}{11pt plus 2pt minus 2pt}
\setlength{\textfloatsep}{12pt plus 2pt minus 2pt}
\setlength{\intextsep}{10pt plus 2pt minus 2pt}
\setlength{\abovedisplayskip}{8pt plus 2pt minus 2pt}
\setlength{\belowdisplayskip}{8pt plus 2pt minus 2pt}
\makeatletter
\renewcommand\section{\@startsection{section}{1}{\z@}%
  {-3.5ex \@plus -1ex \@minus -.2ex}%
  {2.2ex \@plus .25ex}%
  {\normalfont\Large\bfseries}}
\renewcommand\subsection{\@startsection{subsection}{2}{\z@}%
  {-3.1ex\@plus -1ex \@minus -.2ex}%
  {1.5ex \@plus .2ex}%
  {\normalfont\large\bfseries}}
\makeatother

% Contribution list: a bit more room between bullets.
\let\olditemize\itemize
\renewcommand{\itemize}{%
  \olditemize
  \setlength{\itemsep}{0.5ex plus 0.15ex}%
  \setlength{\parsep}{0.15ex}%
}

\title{Trail: A Deterministic-First Log Template Miner with a Per-Decision Audit Trail}
\author{Eduard Beloiu \\ \small MbitAI}
\date{September 2026}

\begin{document}
\maketitle

\begin{abstract}
Log template miners usually return templates and aggregate scores, but
not the sequence of clustering decisions that produced them. Trail is
a deterministic, single-pass miner that writes one audit record per
input line: cluster, similarity, decision, and resulting template. An
optional local 2B-parameter model reviews low-confidence joins without
changing the deterministic parse. On three LogHub-2k datasets, one
global Trail configuration ties tuned Drain on Linux grouping accuracy
(GA 0.690), but trails it on OpenSSH (0.542 vs.\ 0.789); on Apache its
GA is 0.984 vs.\ 1.000. Ablations show that anchors help on Linux and
OpenSSH but fragment Apache, while similarity thresholds from 0.4 to
0.6 do not change these runs. On the synthetic SecOps-2k set, Trail
improves GA from 0.772 to 0.967 and grouping F1 from 0.295 to 0.863
over Drain. A case study shows how line-level records distinguish a
defensible disagreement with ground truth from a low-confidence merge
error. The evaluation uses pinned configurations, dataset hashes, and
reproducible result artifacts.
\end{abstract}

\section{Introduction}

Template mining turns raw log lines into event templates by grouping
lines that differ only in variable parts. Standard evaluations compare
the final groups and templates with labeled data
\cite{jiang2024, zhu2023}. These outputs do not identify the earlier
decision that caused a bad merge. When a parser joins two security
events that should stay apart, an aggregate score measures the effect
but does not record when the join happened or why the parser accepted
it.

Trail records each clustering decision in one JSON record per input
line. Human review and model assist cite those records instead of
re-parsing, preserving a direct path from each final template to the
decisions that formed it. This report reviews related work
(\S\ref{sec:related}), then describes the miner (\S\ref{sec:design}),
the optional local-model review loop (\S\ref{sec:assist}), a worked
case study of two errors hidden by one metric
(\S\ref{sec:case}), a three-parser evaluation on LogHub-2k plus an
ablation (\S\ref{sec:eval}), known limits (\S\ref{sec:limits}), and
reproduction instructions (\S\ref{sec:repro}).

This technical report makes three contributions:
\begin{itemize}
\item a deterministic miner that records the evidence for every
clustering decision;
\item a separate local-model review path whose proposals cannot alter
the base parse without passing explicit gates; and
\item an evaluation with a per-decision OpenSSH case study and
configuration ablations that expose both gains and failure modes.
\end{itemize}

\section{Related work}
\label{sec:related}

Spell and Drain are streaming, syntax-based parsers that cluster log
tokens without model inference \cite{du2016, he2017}. The logparser
toolkit standardized comparisons across multiple datasets and parsers
\cite{zhu2019tools}; LogHub-2.0 later expanded the data and introduced
the GA, PA, FGA, and FTA measures used here \cite{jiang2024}.

Recent systems use language models to extract or refine templates.
EFParser combines small language models with exemplar selection,
caches, and template correction \cite{wang2026efparser}. DeepParse
uses a model once to synthesize regex masks, then applies them through
a deterministic Drain runtime \cite{shetaia2026deepparse}. Matryoshka
instead parses security logs into semantic schemas for downstream
analysis \cite{piet2025security}. Trail does not attempt to match the
accuracy scope of these systems. It focuses on a narrower operational
property: preserving each clustering decision, and keeping optional
model review separate from the deterministic output.

\section{Design}
\label{sec:design}

The miner is one pass, no models, no training. Each line joins the best
matching cluster or starts a new one.

\paragraph{Preprocessing.} The syslog header is stripped with one
regex; only the message part is mined. On LogHub-2k files the harness
splits the header with the same per-dataset format string used for the
Drain baseline, so all parsers see identical input. Messages are split
on whitespace. A small pinned list of whole-token masks replaces
matching tokens with field-named wildcards (e.g.\ \texttt{SRC=}\wild),
matching ground-truth convention.

\paragraph{Candidates.} Only clusters whose token counts differ by at
most \texttt{length\_slack} (1) and whose first \texttt{anchor\_tokens}
(2) match exactly can merge. A candidate is further rejected when
template and line disagree on an \texttt{identity\_keys} field: the
same \texttt{key=} prefix with different values
(\texttt{PROTO=TCP} vs.\ \texttt{PROTO=UDP}). Identity keys name
closed-vocabulary fields that separate templates; everything else may
generalize through key-aware merge.

\paragraph{Join.} Among candidates, the highest token similarity wins.
Exact tokens and any template position holding \wild{} count as
matches. For a template $T$ with $m$ tokens and a line $L$ with $n$
tokens, Trail computes
\begin{equation*}
s(T,L)=\frac{1}{\max(m,n)}
\sum_{i=1}^{\min(m,n)}
\mathbf{1}\bigl[T_i=L_i\ \lor\ T_i=\mbox{\wild}\bigr].
\end{equation*}
The best candidate joins when $s(T,L)\geq\texttt{st}$; the global
configuration in this report sets \texttt{st} to 0.5. Otherwise Trail
opens a new cluster. On join, differing
positions generalize: tokens sharing a \texttt{key=} prefix become
\texttt{key=}\wild{}, anything else becomes bare \wild{}, permanently.
Templates only ever generalize, never split.

\paragraph{Why anchors.} Similarity alone merges lines that differ in
exactly one outcome word. Lines beginning \texttt{Failed password}
and \texttt{Accepted password} agree on every later token, so a
pure threshold puts them in the same template. Requiring the first two
tokens to match keeps outcomes apart without per-dataset tuning. It is
a heuristic, not a principle, and the audit log shows everywhere it
fires, including where it fires wrongly (\S\ref{sec:eval}).

\paragraph{Audit schema (v1).} One record per line:
\begin{lstlisting}
{"line": 12, "cluster": "T3", "decision": "matched",
 "similarity": 0.875, "template": "Failed password for <*> ..."}
\end{lstlisting}
\texttt{decision} is \texttt{new\_cluster} (similarity 0.0) or
\texttt{matched}. Readers must ignore unknown fields; writers must not
rename existing ones. Model proposals and accept/reject decisions live
in a separate append-only review log that cites audit line numbers, so
the parse audit stays an immutable receipt.

\section{Local-model assist}
\label{sec:assist}

The deterministic core stays the parser of record. A local small model
(Qwen3.8-2B distill, Q6\_K, \texttt{empero-ai/Qwen3.8-2B-Distill-GGUF}
\cite{qwen-gguf}, Apache-2.0) reviews candidates selected from the audit log:
low-confidence joins (similarity under 0.7) and near-duplicate cluster
pairs. Each candidate includes three cited example lines. The prompt
(\texttt{trail-lm-v2}) defines SAME vs.\ TWO, adds domain rules
(Accepted vs.\ Failed and BLOCK vs.\ ALLOW are TWO; IP/port/user-only
diffs are SAME; unsure defaults to SAME), and asks one question per
candidate. Anything that is not one unambiguous SAME or TWO is
rejected. Accepted changes land in a separate assisted CSV; the
deterministic CSV and audit never change. Accepts touching more than
10 lines are held as \texttt{needs-human} and never materialized.

The first prompt version (\texttt{trail-lm-v1}, bare SAME/TWO question)
degraded SecOps-2k tight GA from 0.967 to 0.682 through two false
accepts. The v2 preamble plus the hold rule contain exactly that
failure mode: v2 makes zero auto-applies and holds three cases for
human review, so the assisted CSV equals the deterministic parse. We
include the v1 result because it motivates the v2 gate.

Measured cost on an Apple M2 with 8~GB RAM (CPU, llama.cpp server
\cite{llamacpp}):
model load about 40~s at 3.7~GB resident; 16 SecOps-2k candidates in
463~s (about 29~s each); verdicts 13 reject, 3 needs-human, 0
auto-apply. The local run sent no log data to a remote service. This
is one measurement of one model and prompt on one machine, not a
performance benchmark.

\section{Case study: two merges, one number}
\label{sec:case}

On OpenSSH\_2k, Trail trails Drain in grouping accuracy (GA 0.542 vs.\
0.789). Nearly all of the gap sits in two parsed clusters holding 908
of 2000 lines. GA compresses them into one number; the audit records
show they are qualitatively different events.

\paragraph{Exhibit A: confident merge.} Cluster T5 (494 lines) merges
two truth templates that differ by a trailing optional \texttt{user=}
field. Listing~\ref{lst:records} shows the birth record (line~5) and
the merge record (line~28, similarity 0.90); later joins of the same
variant come in at 1.0. The merge is stable, not borderline. A human
with this receipt can reasonably side \emph{with the miner} against
ground truth: same event, optional field. GA counts all 110 lines
wrong either way. Scores cannot express ``wrong by the metric, right
by the event''; the receipt preserves the disagreement in a form a
reviewer can check.

\paragraph{Exhibit B: borderline swallowing.} Cluster T8 (414 lines)
absorbs a single \texttt{disconnected by user} line into 413
\texttt{Bye Bye [preauth]} lines at line~964, similarity 0.625, some
950 lines after the cluster was born. GA barely changes for one line
in 2000, while template-level FGA counts the failed template. Because
the similarity is below 0.7, the LM-assist loop in
\S\ref{sec:assist} selects this record for review. The same evidence
explains the parse and supplies the review queue.

Together, the records show why per-decision evidence matters. The
first merge may reflect a ground-truth boundary that is too strict;
the second is a low-confidence error routed to review by its own
record.

\begin{lstlisting}[basicstyle=\ttfamily\footnotesize,caption={The four audit records behind both merges (line numbers and similarities verbatim).},label={lst:records}]
{"line": 5, "cluster": "T5", "decision": "new_cluster",
 "similarity": 0.0, "template": "pam_unix(sshd:auth):
 authentication failure; logname= uid=0 euid=0 tty=ssh
 ruser= rhost=173.234.31.186 "}
{"line": 28, "cluster": "T5", "decision": "matched",
 "similarity": 0.9, "template": "pam_unix(sshd:auth):
 authentication failure; logname= uid=0 euid=0 tty=ssh
 ruser= rhost=<*> <*>"}
{"line": 14, "cluster": "T8", "decision": "new_cluster",
 "similarity": 0.0, "template": "Received disconnect from
 52.80.34.196: 11: Bye Bye [preauth]"}
{"line": 964, "cluster": "T8", "decision": "matched",
 "similarity": 0.625, "template": "Received disconnect
 from <*> 11: <*> <*> <*>"}
\end{lstlisting}

\section{Evaluation}
\label{sec:eval}

Scorer for all runs: an independent Apache-2.0 implementation of the
LogHub-2.0 metrics \cite{jiang2024}. Grouping accuracy (GA) is the
fraction of messages whose predicted group exactly matches their
ground-truth group. Parsing accuracy (PA) is the fraction whose
predicted template tokens exactly match the ground truth. FGA and FTA
are template-level F1 measures for grouping and template accuracy.
Data:
LogHub-2.0 \texttt{2k\_dataset} at pinned commit
\texttt{ac4aad2} \cite{loghub20}, plus the synthetic SecOps-2k set
(2000 sshd, sudo, firewall, and auditd lines; 25 tight templates).
Each result is one deterministic run on a fixed 2000-line input, so
variance estimates and significance tests do not apply.

\subsection{Tier A: three parsers, three datasets}

Drain uses the logpai benchmark settings \cite{he2017}, which vary by
dataset: \texttt{st} 0.5/0.39/0.6 and depth 4/6/5. Spell runs with an
untuned $\tau=0.5$ \cite{du2016} and the same file formats. Trail runs
one global configuration everywhere: \texttt{st}=0.5, two anchors,
and slack 1. The same Trail configuration produces the SecOps results.

\begin{table}[H]
\centering
\small
\setlength{\tabcolsep}{4pt}
\caption{LogHub-2k results. $N_p$ is the number of predicted templates.}
\label{tab:tiera}
\begin{tabular}{llccccc}
\toprule
Dataset & Parser & GA & PA & FGA & FTA & $N_p$ \\
\midrule
Apache & Drain & 1.000 & 0.694 & 1.000 & 0.500 & 6 \\
Apache & Spell & 0.576 & 0.270 & 0.727 & 0.182 & 5 \\
Apache & Trail & 0.984 & 0.694 & 0.233 & 0.140 & 37 \\
\midrule
Linux & Drain & 0.690 & 0.184 & 0.930 & 0.435 & 112 \\
Linux & Spell & 0.598 & 0.088 & 0.736 & 0.151 & 94 \\
Linux & Trail & 0.690 & 0.171 & 0.918 & 0.395 & 115 \\
\midrule
OpenSSH & Drain & 0.789 & 0.508 & 0.880 & 0.440 & 24 \\
OpenSSH & Spell & 0.252 & 0.121 & 0.488 & 0.195 & 15 \\
OpenSSH & Trail & 0.542 & 0.508 & 0.784 & 0.431 & 25 \\
\bottomrule
\end{tabular}
\end{table}

Trail ties Drain's GA on Linux with one global configuration, although
its PA and template-level scores remain lower. Spell at the default
$\tau$ trails both parsers on all three datasets; tuning Spell could
change that comparison. On Apache, Spell merges the
\texttt{jk2\_init} Found/Can't-find pair that Trail's anchors keep
apart. Trail's Apache and OpenSSH differences are grouping errors
rather than token errors: PA ties Drain on both datasets.

\begin{table}[H]
\centering
\small
\caption{SecOps-2k tight-group results (25 ground-truth templates).}
\label{tab:secops}
\begin{tabular}{lccccc}
\toprule
Parser & GA & PA & FGA & FTA & $N_p$ \\
\midrule
Drain & 0.772 & 0.695 & 0.295 & 0.253 & 70 \\
Trail & 0.967 & 0.967 & 0.863 & 0.863 & 26 \\
\bottomrule
\end{tabular}
\end{table}

Table~\ref{tab:secops} is a controlled synthetic comparison, not
evidence of production security-log performance. The loose 10-group
variant is available in the artifact results; FTA is unchanged by
construction because the loose labels alter grouping, not template
strings.

\subsection{Ablation: anchor sensitivity}

Nine configurations were run on each dataset with the same
preprocessing and scorer. The runner changes only the in-memory Trail
configuration, and the \texttt{base} run reproduces the three Trail
goldens. These ablations are preserved as a runner and raw JSON, but
they are not CI goldens. Table~\ref{tab:ablation} reports GA/FGA with
predicted-template counts in parentheses.

\begin{table}[H]
\centering
\small
\caption{Trail ablation. \texttt{st} variants are identical to base on all
three sets and collapsed to one row.}
\label{tab:ablation}
\begin{tabular}{lccc}
\toprule
Variant & Apache & Linux & OpenSSH \\
\midrule
\texttt{base} & 0.984~/~0.233 (37) & 0.690~/~0.918 (115) & 0.542~/~0.784 (25) \\
\texttt{no-anchors} & 1.000~/~1.000 (6) & 0.186~/~0.457 (44) & 0.279~/~0.622 (19) \\
\texttt{anchors-1} & 1.000~/~1.000 (6) & 0.684~/~0.917 (111) & 0.472~/~0.723 (21) \\
\texttt{st-0.4/0.6} & \multicolumn{3}{c}{identical to \texttt{base}} \\
\texttt{slack-0} & 0.984~/~0.233 (37) & 0.690~/~0.915 (116) & 0.790~/~0.868 (27) \\
\texttt{no-identity} & \multicolumn{3}{c}{identical to \texttt{base}} \\
\texttt{no-masks} & \multicolumn{3}{c}{identical to \texttt{base} (GA/FGA/$N_p$)} \\
\bottomrule
\end{tabular}
\end{table}

Anchors have the largest effect, with opposite signs: they are
essential on Linux (GA 0.690 to 0.186 without) and OpenSSH, but they cause the
Apache over-fragmentation (6 to 37 templates). Anchorless Trail and
Drain produce equivalent message groupings on Apache\_2k, even though
their event identifiers differ. The Apache grouping gap therefore
comes from the anchor rule in this configuration.
One anchor keeps nearly all of the Linux gain; the second helps
OpenSSH. Similarity thresholds from 0.4 to 0.6 produce identical
outputs: the structural filters decide the merges, and the threshold
never binds on these three inputs.
Setting slack to 0 forbids both OpenSSH merges from \S\ref{sec:case}
(they differ in length by 1) and reaches nearly the same GA as Drain
(GA 0.790 vs.\ 0.789), but the two parsers still produce different
partitions. Removing masks leaves grouping unchanged, but lowers
OpenSSH PA from 0.508 to 0.466 and FTA from 0.431 to 0.392. Identity
keys have no effect on these LogHub groups. Choosing
\texttt{no-anchors} for Apache and \texttt{slack-0} for OpenSSH after
seeing the labels would close their GA gaps, but would constitute
per-dataset tuning; we do not report a tuned Trail aggregate.

\section{Limitations}
\label{sec:limits}

Header splitting is syntactic; timestamps, pids, and hostnames are
dropped, not parsed. Masks are a small pinned whole-token list, not a
general field parser. The miner is order-dependent: a different line
order can give different clusters. Inputs are pinned for
reproducibility, but order sensitivity was not measured. Long lines
with many variable tokens can fall below threshold and fragment; the
audit log records the resulting decisions.

SecOps-2k was developed alongside Trail. Its templates exercise
identity keys and masks, so it is a sanity check rather than an
independent benchmark. The data has no organic noise, clock skew, or
interleaving and does not support claims about detector training or
production security logs. The assist results cover one model, prompt,
and machine. LogHub-2k ground-truth labels are disputed: SLGParser
reports significant inaccuracies in the manual labels
\cite{slgparser}, and corrected releases exist
\cite{loghub-corrected}. This report scores against the pinned
originals without relabeling. All evaluations use 2k-line sets.

\section{Reproduction}
\label{sec:repro}

Three repositories form the pipeline. LogParser-Harness evaluates
Drain, Spell, and Trail on LogHub-2k and SecOps-2k
\cite{trail-harness}. LogParser-Dataset contains the synthetic
SecOps-2k set and Drain baseline \cite{trail-dataset}.
LogParser-Trail contains the miner, audit format, and model assist
\cite{trail-miner}. The evaluated
pipeline uses Python 3.12+, \texttt{uv} lockfiles, dataset hashes
checked on fetch, and golden JSON for every row in
Table~\ref{tab:tiera}. The SecOps model outcomes are stored in
\texttt{lm\_scores.json}; the ablation and case-study records are
regenerated by their companion scripts. The case-study script reads
the miner's decision state; a committed OpenSSH \texttt{audit.jsonl}
is not part of the current artifact.

\begin{table}[H]
\centering
\small
\setlength{\tabcolsep}{4pt}
\caption{Pinned artifacts (2026-09-05). Concept DOIs cover all versions.}
\label{tab:artifacts}
\begin{tabular}{@{}llll@{}}
\toprule
Artifact & Commit & Tag & Version DOI \\
\midrule
Harness \cite{trail-harness}
  & \texttt{55d491a} & v0.2.1
  & \doiref{10.5281/zenodo.22341500} \\
Dataset \cite{trail-dataset}
  & \texttt{7ab408d} & v0.1.2
  & \doiref{10.5281/zenodo.22341506} \\
Trail \cite{trail-miner}
  & \texttt{a17e1de} & v0.2.1
  & \doiref{10.5281/zenodo.22341504} \\
\midrule
\multicolumn{4}{@{}l@{}}{Concept DOIs:
  \doiref{10.5281/zenodo.22341499},
  \doiref{10.5281/zenodo.22341505},
  \doiref{10.5281/zenodo.22341503}} \\
\bottomrule
\end{tabular}
\end{table}

\clearpage
\begin{thebibliography}{99}
\small
\raggedright
\setlength{\itemsep}{1.55ex plus 0.4ex}
\setlength{\parsep}{0pt}
\setlength{\parskip}{0pt}
\sloppy
\bibitem{jiang2024} Z.~Jiang et al., ``A large-scale evaluation for log
parsing techniques: How far are we?'' ISSTA, 2024.
\url{https://arxiv.org/abs/2308.10828}

\bibitem{zhu2023} J.~Zhu et al., ``LogHub: A large collection of system
log datasets for AI-driven log analytics.'' ISSRE, 2023.
\url{https://arxiv.org/abs/2008.06448}

\bibitem{he2017} P.~He et al., ``Drain: An online log parsing approach
with fixed depth tree.'' ICWS, 2017.
\url{http://jiemingzhu.github.io/pub/pjhe_icws2017.pdf}

\bibitem{du2016} M.~Du and F.~Li, ``Spell: Streaming parsing of system
event logs.'' ICDM, 2016.
\url{https://www.cs.utah.edu/~lifeifei/papers/spell.pdf}

\bibitem{zhu2019tools} J.~Zhu et al., ``Tools and benchmarks for
automated log parsing.'' ICSE, 2019.
\url{https://arxiv.org/abs/1811.03509}

\bibitem{wang2026efparser} M.~Wang and Y.~Huo, ``Small is beautiful: A
practical and efficient log parsing framework.'' FSE, 2026.
\url{https://doi.org/10.1145/3797145}

\bibitem{shetaia2026deepparse} A.~Shetaia and S.~Kauffman,
``DeepParse: Hybrid log parsing with LLM-synthesized regex masks.''
EASE, 2026. \url{https://arxiv.org/abs/2604.20553}

\bibitem{piet2025security} J.~Piet et al., ``Semantic-aware parsing for
security logs.'' 2025. \url{https://arxiv.org/abs/2506.17512}

\bibitem{qwen-gguf} empero-ai, ``Qwen3.8-2B-Distill-GGUF,'' Hugging
Face (Apache-2.0; weights used for \S\ref{sec:assist}).
\url{https://huggingface.co/empero-ai/Qwen3.8-2B-Distill-GGUF}

\bibitem{llamacpp} G.~Gerganov et al., ``llama.cpp,'' GitHub (MIT).
\url{https://github.com/ggerganov/llama.cpp}

\bibitem{loghub20} LogPAI, ``Loghub-2.0,'' GitHub, commit
\texttt{ac4aad2}.
\url{https://github.com/logpai/Loghub-2.0}

\bibitem{slgparser} Y.~Hu, C.~Wang, L.~Zhao, and A.~Yu, ``SLGParser:
Practical and efficient label-free log parsing using large language
models.'' ICDE, 2026.
\url{https://icde2026.github.io/accepted-papers.html}

\bibitem{loghub-corrected} ``Corrected versions of LogHub, LogHub 2.0,
LoFi, and Hybrid log parsing datasets,'' Zenodo, 2026.
\url{https://zenodo.org/records/20752471}

\bibitem{trail-harness} MbitAI, ``LogParser-Harness,'' v0.2.1
(Apache-2.0).
\url{https://github.com/TMFNK/LogParser-Harness}.
doi:\doiref{10.5281/zenodo.22341500}

\bibitem{trail-dataset} MbitAI, ``LogParser-Dataset'' (SecOps-2k),
v0.1.2 (Apache-2.0).
\url{https://github.com/TMFNK/LogParser-Dataset}.
doi:\doiref{10.5281/zenodo.22341506}

\bibitem{trail-miner} MbitAI, ``LogParser-Trail,'' v0.2.1
(Apache-2.0).
\url{https://github.com/TMFNK/LogParser-Trail}.
doi:\doiref{10.5281/zenodo.22341504}
\end{thebibliography}

\end{document}
